\section{Faserorientierung aus der Spritzgießsimulation}
\label{sec:faserorientierung}

Die Faserorientierung der Spritzgießsimulation wurde nach \secref{sec:fe_modell} über die CONVERSE-Zuordnung auf das FE-Netz übernommen und am Zahn T3 des Sektors so ausgewertet, wie der Forschungsbericht FVA 856 I seine Computertomographie auswertet \cite[Kapitel~6.2 und 6.3]{FVA856-I}. Je Element wird der Orientierungstensor aus dem Faserrahmen und den beiden größten Eigenwerten gebildet, je Zone das volumengewichtete Mittel. Daraus folgen die Eigenwerte $\lambda_{1} \le \lambda_{2} \le \lambda_{3}$, der Anisotropie-Index $AI = 1 - \lambda_{1}/\lambda_{3}$ und die Vorzugsrichtung in der Stirnebene. Die Zonen folgen Bild 17 des Berichts mit dem Randstreifen bis $0{,}2$\,mm unter der Kontur, dem Übergangsstreifen bis $0{,}4$\,mm und darunter dem Kern M über dem Fußkreis und U darunter. \figref{fig:faserorientierung_schnitt} zeigt die Faserrichtung je Element in zwei Stirnschnitten, \figref{fig:faserorientierung_zonen} die Zonen mit der simulierten Vorzugsrichtung.

\begin{figure}[htbp]
    \centering
    \includegraphics[width=\textwidth]{kap5/faserorientierung_schnitt}
    \caption{Faserrichtung je Element in zwei Stirnschnitten des Zahns T3 (Zahnmitte und Stirnseite): Linienrichtung ist die Hauptachse der Orientierung, die Farbe der Orientierungsgrad $a_{11}$. Zonen und Streifen von $0{,}2$\,mm in Anlehnung an \cite[Bild~15 und 17, S.~17 und 19]{FVA856-I}.}
    \label{fig:faserorientierung_schnitt}
\end{figure}

\figref{fig:faserorientierung_schnitt} zeigt, dass die Orientierung über der Zahnbreite nicht gleich ist. In der Zahnmitte liegen die Fasern der Flanken- und Kopfzonen axial, ihre Hauptachse steht fast senkrecht auf der Stirnebene, so dass die Linien dort kurz erscheinen. Im Kern M und unter dem Fußkreis laufen sie radial entlang der Zahnachse. An der Stirnseite folgen die Fasern dagegen in allen Randzonen der Kontur in der Stirnebene, mit langen Linien entlang der Flanke und um den Fuß, und im Kern verwirbeln sie ohne Vorzugsrichtung, unter dem Fußkreis liegen sie tangential. Der Orientierungsgrad $a_{11}$ liegt in beiden Schnitten zwischen $0{,}5$ und $0{,}7$, am höchsten in den Fußzonen und unter dem Fußkreis. Die Zonenwerte der folgenden Absätze sind Mittel über die ganze Breite und mischen beide Bilder.

% TK erledigt (Absatz zu Bild 5.1 ergänzt): Zu wenig oder gar nichts zu Bild 5.1 gesagt, obwohl da schön zu erkennen ist, dass sich die Orientierung über der Zahnbreite unterschiedet.

\begin{figure}[htbp]
    \centering
    \includegraphics[width=0.9\textwidth]{kap5/faserorientierung_zonen}
    \caption{Zonen des Zahns T3 mit der Vorzugsrichtung der Simulation in der Stirnebene, in Anlehnung an \cite[Bild~18, S.~19]{FVA856-I}.}
    \label{fig:faserorientierung_zonen}
\end{figure}

% TK erledigt (Bild kleiner, Beschriftungen entzerrt, Legende ergänzt): Die Grafik gefällt mir nicht. Zu groß auf dem Blatt dargestellt und innerhalb der Grafik überlappen sich einige Beschriftungen. Sieht nicht so qualitätiv hochwertig aus, wie Bild 5.1 und eine Legende fehlt auch.

\tabref{tab:faserorientierung_zonen} stellt den Anisotropie-Index der linken Zonen dem CT-Befund gegenüber. Die rechten Zonen liegen innerhalb von $0{,}08$ bei den linken. Über die ganze Zahnbreite gemittelt sind die simulierten Fußrundungen und unteren Flanken mit $AI$ von $0{,}33$ bis $0{,}35$ weit weniger anisotrop als die CT-Zonen mit $0{,}83$ bis $0{,}86$, die Übergangszonen mit $0{,}20$ bis $0{,}35$ gegen $0{,}60$ bis $0{,}76$. Nahe beieinander liegen der Fußgrund mit $0{,}63$ gegen $0{,}71$, die obere Flanke mit $0{,}62$ gegen $0{,}76$, der Kern M mit $0{,}30$ gegen $0{,}35$ und U mit $0{,}67$ gegen $0{,}70$. Die Elemente selbst sind mit einem mittleren größten Eigenwert von $0{,}55$ bis $0{,}66$ deutlich orientiert. Das Zonenmittel verliert seine Anisotropie durch die Neigung der Faserachse. In beiden Hälften der Zahnbreite kippt sie mit entgegengesetztem Vorzeichen aus der Stirnebene, um $25$ bis $75$° im Drittel an der einen und um $45$ bis $69$° an der anderen Stirnseite, so dass sich die Nebendiagonalglieder aufheben. \figref{fig:faserorientierung_drittel} zeigt deshalb den Index je Drittel der Zahnbreite wie Bild 23 des Berichts. In den Dritteln erreichen Fußrundungen und Flanken $0{,}65$ bis $0{,}80$ gegen $0{,}76$ bis $0{,}86$ im CT, der Fußgrund $0{,}64$ bis $0{,}80$ gegen $0{,}71$ und die Übergangszonen $0{,}47$ bis $0{,}68$ gegen $0{,}60$ bis $0{,}76$. Das mittlere Drittel ist in Fußgrund, Fußrundungen, unteren Flanken und den meisten Übergangszonen weniger anisotrop als die äußeren Drittel. Sein $\lambda_{3}$ liegt um $8$ bis $27$\,\% tiefer. Der Bericht findet an seinem Rad $5$ bis $15$\,\% in den Außenzonen und bis $20$\,\% in den Übergangszonen \cite[Bild~23, S.~24]{FVA856-I}.

\begin{table}[htbp]
    \caption{Anisotropie-Index $AI$ und größter Eigenwert $\lambda_{3}$ der linken Zonen des Zahns T3: Spritzgießsimulation über die ganze Zahnbreite und im mittleren Drittel gegen die CT-Zonen des Forschungsberichts \cite[Bild~17 und 23, S.~19 und 24]{FVA856-I}. Vorzugsrichtung in der Stirnebene ($0$° Zahnachse, positiv nach links), \glqq axial\grqq{} wo die Hauptachse entlang der Zahnbreite liegt.}
    \label{tab:faserorientierung_zonen}
    \centering
    \footnotesize
    \begin{tabular}{L{3.4cm}L{1.3cm}L{1.3cm}L{1.6cm}L{1.3cm}L{2.0cm}L{2.0cm}}
        Zone               & $AI$ Sim. & $AI$ Mitte & $\lambda_{3}$ Sim. & $AI$ CT  & Richtung Sim. & Richtung CT \\
        \hline
        Fußgrund LT        & $0{,}63$  & $0{,}64$   & $0{,}53$           & $0{,}71$ & $-79$°        & $-82$°      \\
        Fußrundung LF      & $0{,}34$  & $0{,}68$   & $0{,}40$           & $0{,}86$ & $-40$°        & $-28$°      \\
        Flanke unten LFU   & $0{,}35$  & $0{,}71$   & $0{,}43$           & $0{,}83$ & axial         & $-13$°      \\
        Flanke oben LFO    & $0{,}62$  & $0{,}75$   & $0{,}50$           & $0{,}76$ & axial         & $-19$°      \\
        Kopf S             & $0{,}41$  & $0{,}35$   & $0{,}45$           & $0{,}55$ & axial         & $-89$°      \\
        Übergang unten LÜU & $0{,}28$  & $0{,}59$   & $0{,}39$           & $0{,}76$ & $-48$°        & $-33$°      \\
        Übergang Mitte LÜM & $0{,}20$  & $0{,}54$   & $0{,}37$           & $0{,}70$ & $-15$°        & $-33$°      \\
        Übergang oben LÜO  & $0{,}25$  & $0{,}52$   & $0{,}37$           & $0{,}60$ & $-47$°        & $-33$°      \\
        Kern M             & $0{,}30$  & $0{,}39$   & $0{,}39$           & $0{,}35$ & $-43$°        & $-85$°      \\
        Kern U             & $0{,}67$  & $0{,}69$   & $0{,}49$           & $0{,}70$ & $89$°         & --          \\
    \end{tabular}
\end{table}

\begin{figure}[htbp]
    \centering
    \includegraphics[width=\textwidth]{kap5/faserorientierung_drittel}
    \caption{Anisotropie-Index je Zone über die ganze Zahnbreite und in den Dritteln der Zahnbreite, in Anlehnung an \cite[Bild~23, S.~24]{FVA856-I}.}
    \label{fig:faserorientierung_drittel}
\end{figure}

In der Richtung stimmen Fußgrund und Kern U mit dem Bericht überein. Die Fasern laufen quer zum Zahn, mit $-79$ und $+82$° gegen $-81{,}6$ und $+84{,}4$°. Die Fußrundungen und unteren Übergangszonen liegen in der Stirnebene bei $-40$ bis $-48$° links und $+50$° rechts, gegen $-27{,}5$ und $+28{,}5$° der Fußrundungen und $\mp 32{,}5$° der Übergangszonen im Bericht, also im selben Sinn, aber $12$ bis $21$° steiler. In den oberen Flanken, am Kopf und in der oberen rechten Übergangszone liegt die Hauptachse dagegen entlang der Zahnbreite. Ihr Anteil in der Stirnebene beträgt nur $0{,}01$ bis $0{,}06$, während das CT dort Fasern in der Ebene entlang der Oberfläche sieht. Der Kern M läuft in der Simulation bei $-43$° statt $-85$°. Die wahrscheinliche Ursache liegt in der Füllung. Der virtuelle Verteiler speist die Stegebene, Kranz und Zähne füllen von der Zahnmitte zu beiden Stirnseiten, so dass die Strömung im dünnen Zahn entlang der Breite läuft. Wie das Rad des Berichts angespritzt wurde, nennt dessen Kapitel 6 nicht. Nach \secref{sec:spritzgiessen} ist für ein solches Rad ein zentraler Schirmanguss anzunehmen, der mit dem virtuellen Verteiler in der Stegebene nicht nachgebildet ist. Welchem Bild das reale Teil folgt, bleibt damit offen. Die Simulation steht gegen das CT eines anderen Rades. Gegen die Tiefe des Randstreifens sind Fußgrund, Fußrundungen, untere Flanken und Kern mit höchstens $0{,}04$ im Index unempfindlich. Obere Flanken und Kopf bewegen sich zwischen $0{,}1$ und $0{,}3$\,mm um bis zu $0{,}13$, die Übergangszonen um bis zu $0{,}5$.

% TK erledigt (Schirmanguss als Annahme genannt, Schliffbild nur im Ausblick): Ich vermute, dass es genauso im Schirmanguss angespritzt wurde, zentral zur Drehachse des Rades wie schon in den Kapiteln vorher beschreiben. Aber das konnte ich eben nicht nachbilden. Bin mir nicht sicher, ob das nochmal an der Stelle erwähnt werden sollte. Aber an sich ja, steht es nicht im Bericht.

Über den Umfang des geschlossenen Rades mit 52 Zähnen streut der Index der Fußrundungen zwischen $0{,}11$ und $0{,}53$ und der oberen Flanken zwischen $0{,}23$ und $0{,}79$, in Wellen von $12$ bis $15$ Zähnen. Die Zonenmittel liegen innerhalb von $0{,}02$ bei denen des Taschenkörpers. Die Streuung ist eine Eigenschaft des Spritzgießergebnisses und nicht des Netzes. Der Zahn 27 des Vollrades entspricht dem Zahn 3 des Sektors. Welcher Zahn in die Rechnung übernommen wird, ist damit eine Eingabe des Modells, vgl. \chapref{chap:zusammenfassung}.

Den Einfluss der Orientierung auf die Beanspruchung zeigt das Schema aus Werkstoffquelle und Richtungsquelle. Die Stufen W3 und W4 wurden mit der Orientierung der CONVERSE-Zuordnung und mit der Orientierung der CT-Zonen des Berichts auf demselben Netz gerechnet. Die CT-Datei trägt je Zone konstante Eigenwerte und ebene Richtungen in den Zonenwinkeln des Berichts, ohne Gradient über die Randschicht und ohne Streuung über die Breite. \figref{fig:orientierungseffekt_16nm_qs} zeigt die relative Änderung an den Normpunkten A bis E bei $16$\,Nm und quasistatischer Karte. Mit den CT-Zonen dreht sich das Ritzel um $5$ bis $13$\,\% weniger und der Zahnkopf verschiebt sich um $3$ bis $14$\,\% weniger, am geringsten an C. Die größte Zugspannung der Fußrundung ändert sich um $-8$ bis $+5$\,\%, die größte Druckspannung steigt um $2$ bis $12$\,\%, an C am stärksten. Der Kontaktdruck wird umverteilt, an A, D und E um bis zu $19$\,\% kleiner, an C meist größer, bei W4 und $16$\,Nm an C jedoch kleiner. Mit den Konstanten der Tabelle 4 des Berichts anstelle der Karte bleibt das Bild, $-7$ bis $-12$\,\% in der Drehung, $-14$ bis $+3$\,\% in der Zugspannung, $+1$ bis $+15$\,\% in der Druckspannung. Die Dehnrate DY1 zeigt den größten, die quasistatische Karte den kleinsten Effekt. Als Maß der Anisotropie einer Karte dient das Verhältnis ihrer Elastizitätsmoduln in und quer zur Fließrichtung der Fasern, das bei $80$\,°C für QS $2{,}17$, für DY1 $2{,}09$ und für DY2 $1{,}85$ beträgt. Die Reihenfolge der Effekte folgt diesem Verhältnis nicht, sondern den übrigen Konstanten der Karten, dem Schubmodul und den Querkontraktionszahlen, und der Lastaufteilung im Kontakt. Die Ursache ist nicht aufgelöst. Dass die Rate den Orientierungseffekt nur um wenige Prozentpunkte verschiebt, obwohl sich die Spannungen der Karten bei gleicher Dehnung deutlich unterscheiden, liegt an der Art des Effekts. Er ist eine relative Größe und vergleicht zwei Orientierungen unter derselben Karte und derselben Kraft, und unter vorgegebener Kraft bestimmt die Steifigkeit die Verformung, nicht die Spannung (\secref{sec:stufen_raten_koerper}). Die Höhe der Kurven wirkt auf die Dehnungen und Festigkeiten des Nachweises, kaum auf die FE-Spannung der linearen Stufen. Der Orientierungseffekt liegt damit in der Größenordnung des Körpereffekts aus \secref{sec:stufen_raten_koerper} und unter dem Sprung zwischen den Werkstoffstufen in \secref{sec:waelzweg}. Für W3 und W4 liegt mit der zugeordneten Orientierung die steife Faserrichtung in oberer Flanke und Kopf entlang der Zahnbreite. Ein reales Teil mit Randschichten in der Stirnebene wie im CT wäre in der Zahnbiegung steifer und fester als das Modell.

% TK erledigt (Legende nennt die Werkstoffstufe der Serie, aufgetragen ist je Stufe die Änderung CT-Zonen gegen Zuordnung; Bildunterschrift präzisiert; Ratenfrage mit den Modulverhältnissen der Karten beantwortet, Ursache offen): Entweder verstehe ich die Grafki nicht, aber ich bin der Meinung, dass die Legende den unterschied zwischen W3 und W4 zeigt und nicht die Unterschiede einer gleichen Werkstoffkarte und Materialmodell, aber dafür CT und Spritzgusssimulation. Bitte überprüfen. Und Warum ist der Effekt bei DY1 stärker als DY2. DY2 hat doch noch eine höhere Dehnrate und müsste damit doch noch steifer und fester sine oder ist da schon wieder ein anderer Effekt vorrangig?

\begin{figure}[htbp]
    \centering
    \includegraphics[width=\textwidth]{kap5/orientierungseffekt_16nm_qs}
    \caption{Orientierungseffekt, relative Änderung jeder Größe, wenn an die Stelle der CONVERSE-Zuordnung die CT-Zonen des Berichts treten, bei gleichem Netz und gleicher Karte an den Normpunkten A bis E. Die Legende nennt die Werkstoffstufe, für die der Tausch gerechnet ist (W3 und W4), $16$\,Nm, quasistatisch. Größtwerte über die Zahnbreite, Marker ohne Linie, weil der Verlauf zwischen den Punkten nicht gerechnet ist.}
    \label{fig:orientierungseffekt_16nm_qs}
\end{figure}

% Herkunft:
% Absatz 1 (Auswertung): ADR-119 Decision 1 bis 3 (architecture_decisions.md ab Zeile 2858); FVA856-I.txt Kapitel 6.2
%   und 6.3 (Bild 15 bis 18 S. 17 bis 19, Bild 23 S. 24, Gl. (1) S. 17 AI); 80_output/fe/faserorientierung/bericht.md
%   Kopf (Zonen, Streifen 0,20/0,40 mm, 13320 Elemente).
% Absatz 2 und Tabelle: bericht.md Abschnitt 1 (Zonen gegen CT) und 3 (Drittel): LT 0,63/0,71, LF 0,34/0,86, LFU
%   0,35/0,83, LFO 0,62/0,76, S 0,41/0,55, LÜU 0,28/0,76, LÜM 0,20/0,70, LÜO 0,25/0,60, M 0,30/0,35, U 0,67/0,70;
%   rechte Zonen RF 0,33, RFU 0,39, RFO 0,70, RT 0,61 (Differenz ≤ 0,08); a11 Mittel 0,55 bis 0,66; Drittel AI
%   mitte: LT 0,64, LF 0,68, LFU 0,71, LFO 0,75, S 0,35, LÜU 0,59, LÜM 0,54, LÜO 0,52, M 0,39, U 0,69; λ3 mitte
%   gegen außen −8 bis −27 % (LT −10, LFU −8, S −14, RFU −11, RF −12, RT −27, LÜM −8, RÜO −21, RÜM −10, RÜU −13,
%   U −10); Neigung z−/z+: LF 31,9/−50,8, LFU 45,8/−50,8, LFO 71,7/−64,5, LÜU 30,1/−42,1, LÜM 41,6/−47,7, RFU
%   45,9/−45,3 → "25 bis 75" und "45 bis 69" (ADR-119 Findings). Richtungen: Abschnitt 2 (LT −79,2, RT 81,6, LF
%   −39,6, RF 49,9, LÜU −48,1, LÜM −15,1, LÜO −46,9, LFU/LFO/S/RFO axial mit Anteil 0,05 bis 0,45, M −42,7, U
%   89,4; Bild 18: −81,6/84,4, −27,5/28,5, ∓32,5, −12,8/−18,5/−89,1, −85,3).
% Absatz 4 (Richtung, Ursache, Empfindlichkeit): ADR-119 Findings und Consequences (virtueller Verteiler 238 Punkte
%   r 15,0 bis 19,1 mm in der Stegebene, Füllung von z = 0 zu beiden Stirnseiten; Band 0,1/0,3 mm: ≤ 0,04, obere
%   Flanken/Kopf ≤ 0,13, Übergangszonen ≤ 0,5; Schliffbild offen, FIB-03).
% Absatz 5 (Vollrad): 80_output/fe/faserorientierung_vollrad/ring/bericht.md (LF 0,12 bis 0,53, LFO 0,23 bis 0,75,
%   RFO 0,24 bis 0,79; Zahn 27 = Zahn 3); ADR-119 Amendment 2026-10-13 (Wellen von 12 bis 15 Zähnen, Zonenmittel
%   innerhalb 0,02 beim Taschenkörper; Wahl des Zahns als Eingabe).
% Absatz 6 (2 × 2-Schema): 80_output/fe/kst_e/vergleich_orientierung_AtoE/analyse.md (Drehung −4,8 bis −12,7 %,
%   Kopfverschiebung −3,0 bis −13,2 %, σ1 −7,6 bis +5,3 %, |σ3| +1,4 bis +12,1 %, Kontaktdruck A/D/E bis −18,8 %,
%   C +4 bis +18,5 %, W4 16 Nm C −14,7/−18,8 %); vergleich_orientierung_tab4_AtoE/analyse.md (W3: Drehung −6,8 bis
%   −12,0, σ1 −14,0 bis +3,3, |σ3| +0,6 bis +14,7 %); ADR-119 Amendment 2026-10-13 (DY1 größter, QS kleinster
%   Effekt; Größenordnung des Körpereffekts, kleiner als die Werkstoffstufe); Konsequenz für W3/W4: ADR-119
%   Consequences, ADR-116 Punkt 47 (b).
% Bildwunsch: kap5/faserorientierung_schnitt (fe/faserorientierung/schnitt_T3.svg), faserorientierung_zonen
%   (zonen_T3.svg, in bilder.yaml), faserorientierung_drittel (drittel_T3.svg), orientierungseffekt_16nm_qs
%   (fe/kst_e/vergleich_orientierung_AtoE/orientierungseffekt_LOAD_16NM_QS.svg); faserorientierung_vergleich
%   (vergleich_T3.svg) steht in bilder.yaml, wird hier durch die Tabelle ersetzt.
% Bibwunsch: FVA856-I vorhanden.
% Nomenklatur: AI, lambda_min/lambda_max (2.2) vorhanden; lambda_1 bis lambda_3 und a_11 (Orientierungsgrad,
%   größter Eigenwert des Elementtensors) prüfen.
% Offen: Schliffbild; Anspritzung des Berichtsrades; die Nachtlauf-Varianten mit Tabelle 4 sind nur für W3 gerechnet.
